\documentclass[11pt]{article}
\usepackage[margin=0.9in]{geometry}
\usepackage{amsmath,amssymb,amsthm,booktabs,mathrsfs}
\usepackage[T1]{fontenc}
\usepackage{lmodern,microtype}
\usepackage[colorlinks=true,linkcolor=blue,urlcolor=blue]{hyperref}
\newtheorem{theorem}{Theorem}[section]
\newtheorem{lemma}[theorem]{Lemma}
\newtheorem{proposition}[theorem]{Proposition}
\newtheorem{corollary}[theorem]{Corollary}
\theoremstyle{remark}
\newtheorem{remark}[theorem]{Remark}
\newcommand{\norm}[1]{\left\|#1\right\|}
\title{Toward a four-dimensional Yang--Mills construction\\
\large Notebook 27: an exact composite blocking step}
\author{Research working document}
\date{17 October 2026}
\begin{document}
\maketitle
\begin{abstract}
Fix one primitive ordinary averaging ratio $\ell$ and take a finite
macro ratio $R=\ell^s$. We construct the local chart and Gaussian
of the actual composite constraint $M_\ell^s$, using a separately
chosen direct-$R$ internal contour gauge. A selected interface bond
has exactly one nonlinear consumer at every intermediate level.
Its flat pivot coefficient is $R^{1-d}$, and elimination preserves
all retained coordinates. A finite rectangular-patch argument gives
coercivity of curvature, contour gauge and the composite constraint;
the preceding classical action bounds then yield a retained Gaussian
gap and exponentially decaying covariance uniformly in the old
refinement. We derive the compact gauge-orbit normalization and the
local Haar, delta, pivot, coupling and Gaussian factors of the
one-shot constrained integral. Two minimum comparisons identify its
physical centered branch. Its pre-gauge representative supplies a
common-ball continuation of an inherited gauge-invariant insertion
without changing integration variables or its cutoff.
The construction is a finite small-field step. Uniform complex
nonvanishing, supported inherited interaction bounds, amplitude
control and the remaining sectors of the existence proof require
their separate estimates.
\end{abstract}

The purpose of this construction is to make the factor $R^{-1}$ of
Notebook 25 available in an actual blocking integral. The primitive
ratio remains fixed while $R=\ell^s$ is chosen large enough for the
desired representative bound. The constraint is defined by its
composition; no identity with a separately defined direct-$R$
ordinary average is needed. Likewise, this one-shot integration is
not identified with $s$ already truncated primitive integrations.
The gauge choice and every factor of the new integral are specified
below. All old-depth uniformity refers to $M=\ell^k$, with $\ell,R$
fixed and compatible periodic volume.

